\section{算术的涌现：在稳定地址空间构造 \texorpdfstring{$(\oplus,\otimes)$}{(oplus, otimes)}}\label{sec:emergent-arithmetic}

本节在 POMmin 下于稳定地址空间内生构造自然数对象以及加法、乘法运算，并证明所得结构与经典半环同构。本文将稳定地址空间取为 Zeckendorf 合法数字的有限支撑部分，即定义 \ref{def:Z-fiber} 中的 $Z$；记号上亦可将其视为 $X_\infty$ 的有限支撑子对象。

\subsection{自然数对象的商幺半群构造}
定义有限支撑数字配置空间
$$
D:=\Bigl\{a=(a_k)_{k\ge 1}: a_k\in\NN,\ \#\{k:a_k\neq 0\}<\infty\Bigr\}.
$$
令 Fibonacci 数列 $(F_k)_{k\ge 1}$ 满足 $F_1=F_2=1$，$F_{k+2}=F_{k+1}+F_k$，定义加权值
$$
\mathrm{Val}(a):=\sum_{k\ge 1} a_k F_{k+1}.
$$
在 $D$ 上引入同余关系 $\equiv$，由基准关系与递推关系生成，使得基向量满足
$$
[e_2]=2[e_1],\qquad [e_{k+2}]=[e_{k+1}]+[e_k]\ \ (k\ge 1),
$$
其中 $e_k$ 表示第 $k$ 位单位向量。

\begin{theorem}[自然数的涌现：商幺半群同构]\label{thm:emergent-N}
映射 $\mathrm{Val}:D/\!\equiv\ \to \NN$ 为交换幺半群同构；并且 Zeckendorf 合法数字空间 $Z$ 为每个同余类提供唯一规范代表：对任意 $a\in D$，存在唯一 $c\in Z$ 使得 $a\equiv c$ 且 $\mathrm{Val}(a)=\mathrm{Val}(c)$。
\par\noindent\textnormal{\textbf{依赖：}定义 \ref{def:Z-fiber}，Zeckendorf 唯一表示定理 \cite{Zeckendorf1972}，以及同余关系 $\equiv$ 的递推生成关系。}
\end{theorem}

\begin{proof}
由 Zeckendorf 表示定理 \cite{Zeckendorf1972}，每个 $N\in\NN$ 存在唯一合法数字 $Z(N)\in Z$ 使得 $\mathrm{Val}(Z(N))=N$，从而 $\mathrm{Val}$ 满射。进一步，在商幺半群 $D/\!\equiv$ 中，递推关系给出 $[e_k]=F_{k+1}[e_1]$，于是任意 $[a]$ 由 $\mathrm{Val}(a)$ 唯一刻画，从而得到同构；规范代表的存在性与唯一性由 Zeckendorf 唯一性给出。
\end{proof}

\subsection{折叠与规范形：加法的涌现}
对 $c,d\in Z$ 定义逐位叠加：令
$$
(c+d)_k:=c_k+d_k\in\{0,1,2\}.
$$
由于叠加通常破坏合法语法，需要通过折叠归一化投影回 $Z$。

\begin{definition}[稳定加法]\label{def:stable-add}
定义
$$
c\oplus d := \Fold_\infty(c+d).
$$
其中 $\Fold_\infty$ 取为定义 \ref{def:fold-extension-domain} 给出的配置空间 $D$ 上的扩展折叠算子（其实现与性质见附录 \ref{app:fold-rewriting}）。等价地，亦可通过值映射定义 $c\oplus d := \mathrm{Val}^{-1}(\mathrm{Val}(c)+\mathrm{Val}(d))$。
\end{definition}

\begin{proposition}[折叠归一化的算法命题]\label{prop:fold-normalize-alg}
对任意 $m\ge 1$ 与任意 $\omega\in\Omega_m$，存在确定性算法在有限步骤内输出唯一正规形 $x=\Fold_m(\omega)\in X_m$，并可附带一条由单步局部重写组成的可核验证书链，供验证器逐步检查（附录 \ref{app:fold-rewriting}）。因此，折叠归一化恢复的是\emph{稳定规范代表}，而非纤维内微态差异。
\end{proposition}

\begin{corollary}[折叠公式]\label{cor:fold-add}
对任意 $c,d\in Z$，有 $c\oplus d=\Fold_\infty(c+d)$，并且 $\mathrm{Val}(c\oplus d)=\mathrm{Val}(c)+\mathrm{Val}(d)$。
\end{corollary}

\begin{theorem}[加法的定义性同构]\label{thm:add-isom}
$(Z,\oplus)$ 为交换幺半群，且
$$
\mathrm{Val}:(Z,\oplus)\to(\NN,+)
$$
为同构。
\end{theorem}

\begin{proof}
由定义 $\mathrm{Val}(c\oplus d)=\mathrm{Val}(c)+\mathrm{Val}(d)$，并由 $\mathrm{Val}$ 的双射性将 $(\NN,+)$ 的交换律、结合律与单位元结构传递到 $(Z,\oplus)$。
\end{proof}

\subsection{乘法的涌现与半环同构}
\begin{definition}[稳定乘法]\label{def:stable-mul}
对 $c,d\in Z$，定义
$$
c\otimes d := \mathrm{Val}^{-1}\bigl(\mathrm{Val}(c)\cdot \mathrm{Val}(d)\bigr)\in Z.
$$
\end{definition}

\begin{theorem}[乘法的定义性同构]\label{thm:mul-isom}
$(Z,\oplus,\otimes)$ 为交换半环，且
$$
\mathrm{Val}:(Z,\oplus,\otimes)\to (\NN,+,\times)
$$
为半环同构。
\end{theorem}

\begin{proof}
由定义 $\mathrm{Val}$ 同时保持加法与乘法，且 $\mathrm{Val}$ 为双射；因此交换律、结合律与分配律从 $(\NN,+,\times)$ 传递。
\end{proof}

\subsection{乘法与加法的内生化：极简原语集}
\begin{proposition}[定义链最小化与可消去性（definitional economy）]\label{prop:definitional-economy-arith}
在本节构造中，可将算术结构的原语集压缩为“折叠正规化 $\Fold_\infty$（及其可核验实现）$+$ 零元 $0_Z$ $+$ 后继 $S(c)=\Fold_\infty(c+e_1)$”。在该原语集上：\par
\noindent\textbf{(i)} 稳定加法 $\oplus$ 可由后继迭代内生（见注 \ref{rem:succ-smaller-primitives}）；\par
\noindent\textbf{(ii)} 稳定乘法 $\otimes$ 不必作为原语：其可由“加法迭代”定义的派生运算 $\odot$ 完全消去（定理 \ref{thm:mul-from-add}）。\par
因此，本节中用 $\mathrm{Val}^{-1}$ 给出的乘法定义（定义 \ref{def:stable-mul}）仅是便于陈述的缩写：半环同构结论的内容在于\emph{内部可定义}的算术运算与经典 $(\NN,+,\times)$ 等价，而非预置经典运算作为原语。
\end{proposition}

\begin{theorem}[乘法可由加法与迭代内生]\label{thm:mul-from-add}
定义
$$
c\odot d := \underbrace{c\oplus c\oplus\cdots\oplus c}_{\mathrm{Val}(d)\text{ 次}},
\qquad
\mathrm{Val}(d)=0 \text{ 时取 } 0_Z,
$$
其中 $0_Z$ 为全零序列。则对任意 $c,d\in Z$，有
$$
c\odot d = c\otimes d.
$$
\end{theorem}

\begin{proof}
由 $\mathrm{Val}$ 保持 $\oplus$ 且双射，得
$$
\mathrm{Val}(c\odot d)=\mathrm{Val}(c)\mathrm{Val}(d)=\mathrm{Val}(c\otimes d),
$$
再由双射性推出相等。
\end{proof}

\begin{remark}[后继算子与更小原语集]\label{rem:succ-smaller-primitives}
同理，可通过后继算子 $S(c)=\Fold_\infty(c+e_1)$ 将 $\oplus$ 内生为后继迭代，从而得到更小的原语集（以“后继 $+$ 折叠”生成算术）。
\end{remark}

